\section{Regularization}
Regularization techniques are methods used to reduce the generalization error of a model, but not necessarily its training error.

Regularization encompasses a variety of strategies designed to mitigate overfitting in machine learning models. 
This section explores the most common approaches applied to Feedforward Neural Networks (FFNNs).

While our primary focus here is on FFNNs, it is worth noting that many of these core methods easily generalize to other architectures, 
such as Convolutional Neural Networks (CNNs), which also rely on their own domain-specific regularization mechanisms.

\subsection{Model Capacity}
Model capacity refers to the ability of a model to fit a wide variety of functions. We have that:

\begin{itemize}
    \item \textbf{Low capacity models} may underfit the data, as they are not able to capture the underlying structure of the data;
    
    \item \textbf{High capacity models} may overfit the data, by memorizing properties of the training data that are not generalizable to unseen data.
\end{itemize}

We can control the capacity of a model by adjusting the hypotheses space, which is the set of all possible functions that the model can represent.
To deal with overfitting, we need to find a good balance between the capacity of the model and its ability to generalize to unseen data.

\subsection{Simple approach}
In addition to controlling the model capacity, many techniques have been proposed to control overfitting, 
some of which where already known in the context of machine learning, such as:

\begin{itemize}
    \item \textbf{Collecting more data}: acquiring more training data can help to reduce overfitting. 
    It allows to better represent the underlying distribution of the data.
    However, this is not always a valiable option, as it might be too expensive or time-consuming to collect more data;
    
    \item \textbf{Ensemble methods}: combining the predictions of multiple models can help to reduce overfitting.
    It allows to average out the errors of individual models, leading to better generalization performance.
    However, this approach is computationally and time-consuming, as it requires training multiple models independently.
\end{itemize}

\subsection{Regularization Methods}
For deeper architectures, controlling the complexity of the model is not as simple as adjusting the number of parameters.
Many techniques have been proposed to regularize such complex models, without reducing their capacity.
Commonly, many regularization methods can be combined to achieve better performance, as they often have complementary effects on the model.

In this section, we will discuss some of the most common regularization methods used in the context of FFNNs.

\subsubsection{Parameter Norm Penalties}

One of the most common regularization techniques is to add a penalty term to the loss function that is propotional to the norm of the model parameters.
This encourages the model to have smaller weights, which has been shown to reduce overfitting.

The general formulation of this regularization technique can be expressed as follows:

\begin{equation*}
    \loss_{\text{reg}}(\vct{\theta}) = \loss(\vct{\theta}) + \lambda \Omega(\vct{\theta})
\end{equation*}

where $\loss(\vct{\theta})$ is the original loss function, $\Omega(\vct{\theta})$ is the regularization term, and $\lambda$
is a hyperparameter that controls the strength of the regularization.

The most common choices for the regularization term $\Omega(\vct{\theta})$ are:
\begin{itemize}
    \item \textbf{L2 Regularization}: also known as Ridge Regression, it adds a penalty term that encourages the model to have smaller weights 
    by adding the squared norm of the weights to the loss function. The regularization term can be expressed as follows:
    \begin{equation*}
        \Omega(\vct{\theta}) = \frac{1}{2} \sum_{i} ||\theta_i^2||
    \end{equation*}

    \item \textbf{L1 Regularization}: also known as Lasso Regression, it adds a penalty term proportional to the absolute value of the weights.
    Differently from L2 regularization, L1 regularization favors the sparsity of the weights, setting the weights of less important features to zero. 
    The regularization term can be expressed as follows:
    \begin{equation*}
        \Omega(\vct{\theta}) = \sum_{i} ||\theta_i||
    \end{equation*}
\end{itemize}

\subsubsection{Data Augmentation}
Data augmentation is a technique used to artificially expand a dataset by applying various transformations to the original samples. 
By exposing the model to a wider variety of inputs, this approach encourages better generalization.
This method is particularly effective for visual data, where operations like rotation, scaling, or flipping can be applied without altering the underlying identity of the object (invariance). 

However, transformations must be carefully tailored to the specific task to avoid corrupting the labels. 
For instance, in the MNIST digit classification dataset, rotating a "6" by 180 degrees transforms it into a "9," inadvertently changing its class.

Furthermore, data augmentation does not guarantee improved performance. 
Because the synthetic examples are merely variations of existing data rather than genuinely novel information, 
they may not always introduce sufficient functional diversity to the model.

\subsubsection{Noise Injection}
Noise injection is a regularization technique that involves adding random noise (variance scaled wrt the model weights) 
to the inputs, outputs, or weights of the model during training.
\begin{itemize}
    \item \textbf{Input Noise}: Adding Gaussian noise to each input example during training can achieve effects similar to $L_2$ regularization. 
    Each example can be represented as:
    \[ x'_i = x_i + \epsilon, \quad \epsilon \sim \mathcal{N}(0, \sigma^2) \]
    where $\mathcal{N}(0, \sigma^2)$ is Gaussian noise with mean 0 and variance $\sigma^2$. 

    For a linear layer, the noise variance is propagated through the weights. 
    If $y = \sum w_i x_i$, the added variance at the node level is proportional to $\sum w_i^2 \sigma^2$. 
    Consequently, to minimize the expected loss, the network is encouraged to maintain smaller weights, mimicking the effect of $L_2$ regularization. 
    High levels of noise typically generate a smoother mapping function, improving generalization performance.

    This concept is also foundational to \textbf{Denoising Autoencoders}, where the network is trained to recover the original input 
    from a corrupted version, forcing it to learn more robust latent representations.

    \item \textbf{Weight Noise}: This technique adds noise to the weights during the forward pass, 
    forcing the model to remain performant even when parameters are slightly perturbed. 
    This encourages the training process to converge to a \textbf{flat minimum} of the loss function, 
    where the model's performance is less sensitive to small changes in weight values.

    \item \textbf{Output Noise}: This discourages the model from making overconfident predictions by smoothing out the label distribution (often referred to as label smoothing). 
    This prevents the network from outputting extremely high logits for a single class, which can lead to better calibration and reduced overfitting.
\end{itemize}

\subsubsection{Early Stopping}
A simple yet effective regularization technique is to stop the training process before the model starts to overfit the training data.
This can be done by monitoring the performance of the model on a validation set during training, and stopping the training process when the 
performance on the validation set starts to degrade.

This procedure can be summarized as follows:
\begin{enumerate}
    \item Start the training process with small weights;
    
    \item Every time the error on the validation set decreases, save a copy of the model parameters;
    
    \item When the training process is stopped, restore the model parameters to the values that achieved the best performance on the validation set.
\end{enumerate}

\subsubsection{Multi-Task Learning}
Multi-task learning is a regularization technique that involves training a model on multiple related tasks simultaneously.
\begin{figure}[H]
    \centering
    % STYLES
\tikzset{
  >=latex, % for default LaTeX arrow head
  node/.style={thick,circle,draw=neuronedge,minimum size=22,inner sep=0.5,outer sep=0.6},
  node in/.style={node,fill=neuronfill},
}
\def\nstyle{int(\lay<\Nnodlen?min(2,\lay):3)} % map layer number onto 1, 2, or 3

\begin{tikzpicture}[x=2.7cm,y=1.6cm]
  \draw[->] (0.14, 0) -- (0.50, 0) node[pos=0.50] {};
  \node[node in] (0, 0) at (0, 0) {$x$};
  \draw[->] (0.84, 0) -- (1.26, 0) node[pos=0.50] {};
  \draw[->] (0.78, 0.3) -- (1.26, 1) node[pos=0.50] {};  
  \draw[->] (0.78, -0.3) -- (1.26, -1) node[pos=0.50] {};  
  \node[node in] (0, 0) at (0.7, 0) {$h^\text{shared}$};
  \draw[->] (1.54, 1) -- (1.96, 1) node[pos=0.50] {};
  \node[node in] (0, 0) at (1.4, 1) {$h^{(1)}$};
  \draw[->] (1.54, 0) -- (1.96, 0) node[pos=0.50] {};
  \node[node in] (0, 0) at (1.4, 0) {$h^{(2)}$};
  \node[node in] (0, 0) at (1.4, -1) {$h^{(3)}$};
  \node[node in] (0, 0) at (2.1, 1) {$y^{(1)}$};
  \node[node in] (0, 0) at (2.1, 0) {$y^{(2)}$};
  
\end{tikzpicture}

    \caption{Example of a multi-task learning architecture. 
    The model is trained on multiple related tasks simultaneously, sharing some intermediate representations.}
    \label{fig:multi-task-learning}
\end{figure}

Usually, we have several tasks that share the same set of input data $\vct{x}$, as well as some intermediate representations $h^{(shared)}$,
but may have different output layers $f^{(1)}$, $f^{(2)}$, ..., $f^{(n)}$ for each task.

By sharing the intermediate representations, the model has 2 advantages:
\begin{itemize}
    \item It can learn more robust and generalizable features, as it is forced to learn representations that are useful for multiple tasks;
    \item It can leverage the additional data from the other tasks, which can help to reduce overfitting and improve generalization performance.
\end{itemize}

\subsubsection{Parameter Sharing}
Parameter sharing (or parameter tying) is a regularization strategy that reduces model complexity by forcing different parts of the network to use the same weights. 
This can be achieved softly, by adding a penalty to the loss function that encourages weight similarity, or explicitly, by tying the parameters together during optimization so they remain identical.

This approach is particularly powerful when prior knowledge about the data structure suggests that certain patterns are universal. 
The prime example is the \textbf{Convolutional Neural Network} (CNN), which applies a single, shared set of weights (a filter) across all spatial locations of an input. 
This structural constraint dramatically reduces the parameter count and allows the model to efficiently learn translation-invariant features.


\subsubsection{Dropout}
We have already seen that the ensemble of multiple models can help to reduce overfitting by averaging out the errors of individual models.
In deep learning, this approach would be computationally and time-consuming, as it would require training multiple models independently.

\textbf{Dropout} is a regularization technique born from the idea of ensemble learning, which allows to approximate the effect of training 
multiple models by randomly dropping out a subset of the neurons during training.
\begin{figure}[H]
    \centering
    \begin{tikzpicture}

	\node[circle, draw, thick, draw=neuronedge, fill=neuronfill] (i1) {};
	\node[circle, draw, thick, draw=neuronedge, fill=neuronfill, above=2em of i1] (i2) {};
	\node[circle, draw, thick, draw=neuronedge, fill=neuronfill, above=2em of i2] (i3) {};
	\node[circle, draw, thick, draw=neuronedge, fill=neuronfill, below=2em of i1] (i4) {};
	\node[circle, draw, thick, draw=neuronedge, fill=neuronfill, below=2em of i4] (i5) {};
	
	\node[circle, draw, thick, draw=neuronedge, fill=neuronfill, right=4em of i1] (h1) {};
	\node[circle, draw, thick, draw=neuronedge, fill=neuronfill, right=4em of i2] (h2) {};
	\node[circle, draw, thick, draw=neuronedge, fill=neuronfill, right=4em of i3] (h3) {};
	\node[circle, draw, thick, draw=neuronedge, fill=neuronfill, right=4em of i4] (h4) {};
	\node[circle, draw, thick, draw=neuronedge, fill=neuronfill, right=4em of i5] (h5) {};
	
	\node[circle, draw, thick, draw=neuronedge, fill=neuronfill, right=4em of h1] (hh1) {};
	\node[circle, draw, thick, draw=neuronedge, fill=neuronfill, right=4em of h2] (hh2) {};
	\node[circle, draw, thick, draw=neuronedge, fill=neuronfill, right=4em of h3] (hh3) {};
	\node[circle, draw, thick, draw=neuronedge, fill=neuronfill, right=4em of h4] (hh4) {};
	\node[circle, draw, thick, draw=neuronedge, fill=neuronfill, right=4em of h5] (hh5) {};
	
	\node[circle, draw, thick, draw=neuronedge, fill=neuronfill, right=4em of hh2] (o1) {};
	\node[circle, draw, thick, draw=neuronedge, fill=neuronfill, right=4em of hh4] (o2) {};
	
	\draw[-, thick, thick,conncolor] (i1) -- (h1);
	\draw[-, thick, thick,conncolor] (i1) -- (h2);
	\draw[-, thick, thick,conncolor] (i1) -- (h3);
	\draw[-, thick, thick,conncolor] (i1) -- (h4);
	\draw[-, thick, thick,conncolor] (i1) -- (h5);
	\draw[-, thick, thick,conncolor] (i2) -- (h1);
	\draw[-, thick, thick,conncolor] (i2) -- (h2);
	\draw[-, thick, thick,conncolor] (i2) -- (h3);
	\draw[-, thick, thick,conncolor] (i2) -- (h4);
	\draw[-, thick, thick,conncolor] (i2) -- (h5);
	\draw[-, thick, thick,conncolor] (i3) -- (h1);
	\draw[-, thick, thick,conncolor] (i3) -- (h2);
	\draw[-, thick, thick,conncolor] (i3) -- (h3);
	\draw[-, thick, thick,conncolor] (i3) -- (h4);
	\draw[-, thick, thick,conncolor] (i3) -- (h5);
	\draw[-, thick, thick,conncolor] (i4) -- (h1);
	\draw[-, thick, thick,conncolor] (i4) -- (h2);
	\draw[-, thick, thick,conncolor] (i4) -- (h3);
	\draw[-, thick, thick,conncolor] (i4) -- (h4);
	\draw[-, thick, thick,conncolor] (i4) -- (h5);
	\draw[-, thick, thick,conncolor] (i5) -- (h1);
	\draw[-, thick, thick,conncolor] (i5) -- (h2);
	\draw[-, thick, thick,conncolor] (i5) -- (h3);
	\draw[-, thick, thick,conncolor] (i5) -- (h4);
	\draw[-, thick, thick,conncolor] (i5) -- (h5);
	
	\draw[-, thick, thick,conncolor] (h1) -- (hh1);
	\draw[-, thick, thick,conncolor] (h1) -- (hh2);
	\draw[-, thick, thick,conncolor] (h1) -- (hh3);
	\draw[-, thick, thick,conncolor] (h1) -- (hh4);
	\draw[-, thick, thick,conncolor] (h1) -- (hh5);
	\draw[-, thick, thick,conncolor] (h2) -- (hh1);
	\draw[-, thick, thick,conncolor] (h2) -- (hh2);
	\draw[-, thick, thick,conncolor] (h2) -- (hh3);
	\draw[-, thick, thick,conncolor] (h2) -- (hh4);
	\draw[-, thick, thick,conncolor] (h2) -- (hh5);
	\draw[-, thick, thick,conncolor] (h3) -- (hh1);
	\draw[-, thick, thick,conncolor] (h3) -- (hh2);
	\draw[-, thick, thick,conncolor] (h3) -- (hh3);
	\draw[-, thick, thick,conncolor] (h3) -- (hh4);
	\draw[-, thick, thick,conncolor] (h3) -- (hh5);
	\draw[-, thick, thick,conncolor] (h4) -- (hh1);
	\draw[-, thick, thick,conncolor] (h4) -- (hh2);
	\draw[-, thick, thick,conncolor] (h4) -- (hh3);
	\draw[-, thick, thick,conncolor] (h4) -- (hh4);
	\draw[-, thick, thick,conncolor] (h4) -- (hh5);
	\draw[-, thick, thick,conncolor] (h5) -- (hh1);
	\draw[-, thick, thick,conncolor] (h5) -- (hh2);
	\draw[-, thick, thick,conncolor] (h5) -- (hh3);
	\draw[-, thick, thick,conncolor] (h5) -- (hh4);
	\draw[-, thick, thick,conncolor] (h5) -- (hh5);
	
	
	\draw[-, thick, thick,conncolor] (hh1) -- (o1);
	\draw[-, thick, thick,conncolor] (hh1) -- (o2);
	\draw[-, thick, thick,conncolor] (hh2) -- (o1);
	\draw[-, thick, thick,conncolor] (hh2) -- (o2);
	\draw[-, thick, thick,conncolor] (hh3) -- (o1);
	\draw[-, thick, thick,conncolor] (hh3) -- (o2);
	\draw[-, thick, thick,conncolor] (hh4) -- (o1);
	\draw[-, thick, thick,conncolor] (hh4) -- (o2);
	\draw[-, thick, thick,conncolor] (hh5) -- (o1);
	\draw[-, thick, thick,conncolor] (hh5) -- (o2);
	
	\draw[-stealth, very thick] (5.5,0) -- node[above] {} (7, 0);
	
	
	\node[circle, draw, thick, draw=neuronedge, fill=neuronfill, draw=dlred, fill=dlred!12, right=10.5em of hh1] (i1) {};
	\node[circle, draw, thick, draw=neuronedge, fill=neuronfill, draw=dlred, fill=dlred!12, above=2em of i1] (i2) {};
	\node[circle, draw, thick, draw=neuronedge, fill=neuronfill, above=2em of i2] (i3) {};
	\node[circle, draw, thick, draw=neuronedge, fill=neuronfill, below=2em of i1] (i4) {};
	\node[circle, draw, thick, draw=neuronedge, fill=neuronfill, below=2em of i4] (i5) {};
	
	\node[dlred] (icr) at (i1) {$\mathlarger{\mathlarger{\mathlarger{\mathlarger{\mathlarger{\bm{\times}}}}}}$};
	\node[dlred] (icr) at (i2) {$\mathlarger{\mathlarger{\mathlarger{\mathlarger{\mathlarger{\bm{\times}}}}}}$};
	
	\node[circle, draw, thick, draw=neuronedge, fill=neuronfill, draw=dlred, fill=dlred!12, right=4em of i1] (h1) {};
	\node[circle, draw, thick, draw=neuronedge, fill=neuronfill, right=4em of i2] (h2) {};
	\node[circle, draw, thick, draw=neuronedge, fill=neuronfill, draw=dlred, fill=dlred!12, right=4em of i3] (h3) {};
	\node[circle, draw, thick, draw=neuronedge, fill=neuronfill, draw=dlred, fill=dlred!12, right=4em of i4] (h4) {};
	\node[circle, draw, thick, draw=neuronedge, fill=neuronfill, right=4em of i5] (h5) {};
	
	\node[dlred] (icr) at (h1) {$\mathlarger{\mathlarger{\mathlarger{\mathlarger{\mathlarger{\bm{\times}}}}}}$};
	\node[dlred] (icr) at (h3) {$\mathlarger{\mathlarger{\mathlarger{\mathlarger{\mathlarger{\bm{\times}}}}}}$};
	\node[dlred] (icr) at (h4) {$\mathlarger{\mathlarger{\mathlarger{\mathlarger{\mathlarger{\bm{\times}}}}}}$};
	
	\node[circle, draw, thick, draw=neuronedge, fill=neuronfill, right=4em of h1] (hh1) {};
	\node[circle, draw, thick, draw=neuronedge, fill=neuronfill, draw=dlred, fill=dlred!12, right=4em of h2] (hh2) {};
	\node[circle, draw, thick, draw=neuronedge, fill=neuronfill, right=4em of h3] (hh3) {};
	\node[circle, draw, thick, draw=neuronedge, fill=neuronfill, draw=dlred, fill=dlred!12, right=4em of h4] (hh4) {};
	\node[circle, draw, thick, draw=neuronedge, fill=neuronfill, right=4em of h5] (hh5) {};
	
	\node[dlred] (icr) at (hh2) {$\mathlarger{\mathlarger{\mathlarger{\mathlarger{\mathlarger{\bm{\times}}}}}}$};
	\node[dlred] (icr) at (hh4) {$\mathlarger{\mathlarger{\mathlarger{\mathlarger{\mathlarger{\bm{\times}}}}}}$};
	
	\node[circle, draw, thick, draw=neuronedge, fill=neuronfill, right=4em of hh2] (o1) {};
	\node[circle, draw, thick, draw=neuronedge, fill=neuronfill, right=4em of hh4] (o2) {};
	
	\draw[-, thick, thick,conncolor] (i3) -- (h2);
	\draw[-, thick, thick,conncolor] (i3) -- (h5);
	\draw[-, thick, thick,conncolor] (i4) -- (h2);
	\draw[-, thick, thick,conncolor] (i4) -- (h5);
	\draw[-, thick, thick,conncolor] (i5) -- (h2);
	\draw[-, thick, thick,conncolor] (i5) -- (h5);
	
	\draw[-, thick, thick,conncolor] (h2) -- (hh1);
	\draw[-, thick, thick,conncolor] (h2) -- (hh3);
	\draw[-, thick, thick,conncolor] (h2) -- (hh5);
	\draw[-, thick, thick,conncolor] (h5) -- (hh1);
	\draw[-, thick, thick,conncolor] (h5) -- (hh3);
	\draw[-, thick, thick,conncolor] (h5) -- (hh5);
	
	\draw[-, thick, thick,conncolor] (hh1) -- (o1);
	\draw[-, thick, thick,conncolor] (hh1) -- (o2);
	\draw[-, thick, thick,conncolor] (hh3) -- (o1);
	\draw[-, thick, thick,conncolor] (hh3) -- (o2);
	\draw[-, thick, thick,conncolor] (hh5) -- (o1);
	\draw[-, thick, thick,conncolor] (hh5) -- (o2);

\end{tikzpicture}
    \caption{Example of a dropout regularization technique.}
    \label{fig:dropout}
\end{figure}
During the training process, each neuron is randomly dropped out (set to zero) with a certain probability $p$, which is a hyperparameter that controls the strength of the regularization.
This forces the model to learn redundant representations, which allows the model to be more robust and less sensitive to the specific weights of individual neurons, 
thus improving generalization performance.

Dropout is similar to a \textit{bagging} approach, where we train multiple models on different subsets of the data and average their predictions.
We can see each trainging iteration as training a different model. At test time, we use the full network, but we scale the weights of the neurons by the dropout probability $p$ to account for the fact that during training, each neuron was only active with probability $p$.

\subsubsection{Implicit Regularization}
Some techinques, such as Stochastic Gradient Descent (SGD), have been shown to have an implicit regularization effect on the model.
In particular, its stochastic nature can favor flatter minima of the loss function, which are associated with better generalization performance.

Also, using loss functions that are less sensitive to outliers, such as the Huber loss, can have an implicit regularization effect by reducing the influence of outliers on the training process.